\documentclass[11pt,a4paper]{article}
\usepackage[margin=2.5cm]{geometry}
\usepackage[T1]{fontenc}
\usepackage{booktabs}
\usepackage{xcolor}
\usepackage{listings}
\usepackage[hidelinks]{hyperref}

\emergencystretch=2em
\definecolor{codebg}{gray}{0.95}
\lstset{language=C++, basicstyle=\ttfamily\small, keywordstyle=\color{blue!70!black},
  commentstyle=\color{green!50!black}\itshape, backgroundcolor=\color{codebg},
  frame=single, rulecolor=\color{black!20}, breaklines=true, showstringspaces=false,
  tabsize=2, columns=flexible}

\title{ATmega32A Driver Course \\ \Large Module 1 --- GPIO (Digital Input / Output)}
\author{Ali Sahafi}
\date{}

\begin{document}
\maketitle

\section{Theory}

GPIO (\emph{General Purpose Input/Output}) is the most basic peripheral: each
pin of the microcontroller can be a digital \textbf{output} (drive 0\,V or
5\,V) or a digital \textbf{input} (read whether the outside world applies
0\,V or 5\,V). The ATmega32A has four 8-bit ports: \textbf{PORTA},
\textbf{PORTB}, \textbf{PORTC} and \textbf{PORTD} --- 32 pins in total.

Each port is controlled by three registers ($x$ = A, B, C or D):

\begin{center}
\begin{tabular}{lll}
\toprule
\textbf{Register} & \textbf{Role} & \textbf{Bit meaning} \\
\midrule
\texttt{DDR$x$}  & Data Direction & 1 = output, 0 = input \\
\texttt{PORT$x$} & Output latch / pull-up & output: pin level; input: 1 = pull-up on \\
\texttt{PIN$x$}  & Input read & the actual voltage currently on the pin \\
\bottomrule
\end{tabular}
\end{center}

\paragraph{Internal pull-up.} An unconnected input pin \emph{floats}: it reads
random values. For buttons, enable the internal pull-up
(\texttt{INPUT\_PULLUP}): the pin then rests at \texttt{HIGH}, and the button
connects it to GND, so \textbf{pressed = LOW}. This ``active-low'' logic is
the standard way to wire buttons.

\section{API reference}

The driver auto-corrects a common mistake: if you accidentally pass
\texttt{PORTx} where \texttt{DDRx} is expected, \texttt{setDirection} fixes
it for you.

\begin{center}
\small
\begin{tabular}{ll}
\toprule
\textbf{Method} & \textbf{Description} \\
\midrule
\texttt{GPIO.setDirection(DDRx, pin, OUTPUT)} & Set single pin direction \\
\texttt{GPIO.setDirection(DDRx, ALL, INPUT\_PULLUP)} & Set whole port direction \\
\texttt{GPIO.write(PORTx, pin, HIGH/LOW)} & Write single pin \\
\texttt{GPIO.write(PORTx, ALL, 0xF0)} & Write raw byte to whole port \\
\texttt{GPIO.write(PORTx, 0xF0)} & Shorthand whole-port write \\
\texttt{GPIO.read(PINx, pin)} & Read single pin --- returns \texttt{HIGH} or \texttt{LOW} \\
\texttt{GPIO.read(PINx, ALL)} & Read whole port --- returns 0--255 \\
\texttt{GPIO.read(PINx)} & Shorthand whole-port read \\
\texttt{GPIO.toggle(PORTx, pin)} & Toggle single pin \\
\texttt{GPIO.toggle(PORTx, ALL)} & Toggle whole port \\
\bottomrule
\end{tabular}
\end{center}

\textbf{Constants:} \texttt{INPUT}, \texttt{OUTPUT}, \texttt{INPUT\_PULLUP},
\texttt{HIGH}, \texttt{LOW}, \texttt{ALL} (see Module 0).

\section{Example: button-controlled LED}

\paragraph{Wiring.}
\begin{itemize}
  \item Push button between \textbf{PD2} and \textbf{GND} (internal pull-up used, no external resistor needed).
  \item LED + 330\,$\Omega$ resistor from \textbf{PC0} to GND (follows the button).
  \item LED + 330\,$\Omega$ resistor from \textbf{PC1} to GND (blinks continuously).
\end{itemize}

\paragraph{Build \& flash.} \texttt{make gpio}

\lstinputlisting{example.cpp}

Note that reading uses \texttt{PIND} (the input register) while writing uses
\texttt{PORTC} (the output latch) --- mixing these up is the classic GPIO bug.

\section{Exercises}

\begin{enumerate}
  \item Change the program so the LED on PC0 \emph{toggles} once per button
        press instead of following it. (Hint: you must detect the
        \emph{edge} --- remember the previous button state in a variable.)
  \item Connect 8 LEDs to PORTC and write a ``running light'' (Knight Rider)
        pattern using the whole-port \texttt{GPIO.write(PORTC, value)} overload.
  \item Why does a button need a pull-up at all? Remove
        \texttt{INPUT\_PULLUP} (use plain \texttt{INPUT}), run the program and
        explain what you observe.
\end{enumerate}

\end{document}
